\documentclass[11pt]{article}

\usepackage[a4paper,margin=0.8in]{geometry}
\usepackage{amsmath}
\usepackage{amssymb}
\usepackage{array}
\usepackage{booktabs}
\usepackage{graphicx}
\usepackage{hyperref}
\usepackage{tabularx}

\title{Topology Comparision}
\author{Chi Zhang and Codex}
\date{\today}

\newcommand{\toponame}[1]{\texttt{#1}}
\newcommand{\Gbps}{\mathrm{GB/s}}
\newcommand{\floor}[1]{\left\lfloor #1 \right\rfloor}
\newcommand{\ceil}[1]{\left\lceil #1 \right\rceil}

\begin{document}
\maketitle

\section*{Metric Conventions}

This report compares the graph models currently used by Topology Analyzer.
Unless otherwise stated, $N$ is the number of terminal or injection nodes,
$R$ is the number of routers, $k$ is the maximum router radix including local
terminal ports, $E$ is the number of bidirectional physical router-router link
pairs, $D$ is router-hop diameter, and $B$ is one-way bisection bandwidth. The
concrete examples assume homogeneous bidirectional links with latency
$1$ cycle and one-way bandwidth $b=1\Gbps$.

The repository's benchmark \texttt{metrics.txt} reports directed
router-router channels for \texttt{links}. For bidirectional examples in this
report, that directed count is $2E$.

\section{Introduction to Topologies Under Analysis}

\subsection{\toponame{mesh2d}}

A 2D mesh arranges routers on a rectangular $x \times y$ grid. Each router is
connected to its immediate horizontal and vertical neighbors when such
neighbors exist. Interior routers therefore have four network links, edge
routers have fewer, and each router attaches $c$ local terminal nodes.

The parameters are $x$ routers in the X dimension, $y$ routers in the Y
dimension, and concentration $c$. The common baseline routing is dimension
order routing, for example XY routing, where packets finish one coordinate
dimension before moving in the next dimension.

\subsection{\toponame{mesh3d}}

A 3D mesh extends the rectangular mesh to $x \times y \times z$ routers.
Routers connect to immediate neighbors in the X, Y, and Z dimensions. This
uses the same local-grid structure as \toponame{mesh2d}, but adds vertical
connectivity, reducing diameter for a fixed node count when the third dimension
is physically practical.

The parameters are $x$, $y$, and $z$ routers in each dimension, plus
concentration $c$.

\subsection{\toponame{dragonfly}}

Dragonfly is a low-diameter grouped topology designed for large-scale high
performance networks \cite{kim2008dragonfly}. Routers are partitioned into
groups. Inside each group, the $a$ routers form a local all-to-all subgraph.
Across groups, global links connect groups so that packets can usually move
from a source group to a destination group in one global hop, with local hops
at the source and destination groups.

The parameters are $p$ terminal nodes per router, $a$ routers per group, $h$
global links per router, and $g$ groups. In this repository, when $g$ is not
specified, the default full arrangement is $g=ah+1$, giving one global link
between every pair of groups.

\subsection{\toponame{hypercube}}

A binary hypercube labels routers with $d$-bit binary coordinates. Two routers
are connected if their labels differ in exactly one bit. This gives exactly
$d$ network links per router, logarithmic diameter, and high path diversity.

The parameters are dimension $d$ and concentration $c$. With $c=1$, the number
of terminal nodes equals the number of routers, $2^d$.

\subsection{\toponame{slimnoc}}

SlimNoC is an on-chip network topology derived from Slim Fly and the
McKay-Miller-Siran diameter-2 graph family \cite{besta2020slimnoc}. Routers
are generated from a finite-field construction. The graph has two router
partitions, each partition is organized into $q$ subgroups of $q$ routers, and
cross links connect the two partitions according to finite-field arithmetic.
The result is a diameter-2 topology with much lower diameter than a mesh at
the cost of higher router radix and less geometric locality.

The main topology parameter is finite-field order $q$, with
$q=4w+\delta$ and $\delta \in \{-1,0,1\}$ for the supported construction.
The concentration $p$ gives local terminal nodes per router. In the repository
YAML this concentration can be written as \texttt{concentration} or \texttt{p}.

\subsection{\toponame{ubmesh}}

UBMesh is modeled here as the nD-FullMesh core described by the UB-Mesh paper
\cite{liao2025ubmesh}. Routers have coordinates
$(u_0,\ldots,u_{d-1})$ in dimensions $[L_0,\ldots,L_{d-1}]$. Along each
dimension, all routers that share the other coordinates are fully connected.
This gives very small diameter, equal to the number of non-singleton
dimensions, while the router radix grows with the sum of dimension lengths.

The parameters are the dimension lengths $L_i$, optional dimension names, and
concentration $c$.

\subsection{\toponame{fattree}}

Fat-tree networks increase path diversity and bisection bandwidth by using a
hierarchical tree-like structure with multiple parents instead of a single
thin root bottleneck \cite{leiserson1985fattree}. The model used here is a
radix-split $k$-ary $n$-tree. For router radix $r$, each router has
$s=r/2$ downward ports and $s=r/2$ upward ports, except at the edges of the
hierarchy where ports attach to terminals or to adjacent tree levels.

The parameters are router radix $r$ and level count $L$. The split is
$s=r/2$. Leaf routers each attach $s$ terminal nodes.

\section{Scaling Laws}

Table~\ref{tab:scaling-laws} summarizes the formulas used for the comparison.
For heterogeneous link bandwidths, replace $b$ with the corresponding
dimension or class bandwidth. The formulas below use the homogeneous case for
compactness.

\begin{table}[htbp]
\centering
\scriptsize
\caption{Scaling laws by topology. $E$ counts bidirectional router-router link pairs. $B$ is one-way bisection bandwidth.}
\label{tab:scaling-laws}
\resizebox{\textwidth}{!}{%
\begin{tabular}{@{}p{0.13\textwidth}p{0.15\textwidth}p{0.16\textwidth}p{0.16\textwidth}p{0.19\textwidth}p{0.13\textwidth}p{0.22\textwidth}@{}}
\toprule
Topology & Parameters & Nodes $N$ & Max router radix $k$ & Links $E$ & Diameter $D$ & Bisection bandwidth $B$ \\
\midrule
\toponame{mesh2d}
& $x,y,c$
& $cxy$
& $c+\min(2,x-1)+\min(2,y-1)$
& $(x-1)y+x(y-1)$
& $(x-1)+(y-1)$
& $\min(y,x)b$ \\
\toponame{mesh3d}
& $x,y,z,c$
& $cxyz$
& $c+\min(2,x-1)+\min(2,y-1)+\min(2,z-1)$
& $(x-1)yz+x(y-1)z+xy(z-1)$
& $(x-1)+(y-1)+(z-1)$
& $\min(yz,xz,xy)b$ \\
\toponame{dragonfly}
& $p,a,h,g$
& $pag$
& $p+(a-1)+\min(h,g-1)$
& $g\binom{a}{2}+\binom{g}{2}$
& $3$ for full cross-group case
& $\floor{g/2}\ceil{g/2}b$ \\
\toponame{hypercube}
& $d,c$
& $c2^d$
& $c+d$
& $d2^{d-1}$
& $d$
& $2^{d-1}b$ \\
\toponame{slimnoc}
& $q,p,\delta$
& $2q^2p$
& $p+\frac{3q-\delta}{2}$
& $q^2\frac{3q-\delta}{2}$
& $2$
& $2(q-1)\floor{q/2}\ceil{q/2}b$ \\
\toponame{ubmesh}
& $[L_0,\ldots,L_{d-1}],c$
& $c\prod_i L_i$
& $c+\sum_i(L_i-1)$
& $\frac{1}{2}\left(\prod_iL_i\right)\sum_i(L_i-1)$
& $\sum_i \mathbf{1}[L_i>1]$
& $\min_i\left(\frac{R}{L_i}\floor{L_i/2}\ceil{L_i/2}b\right)$ \\
\toponame{fattree}
& $r,L,s=r/2$
& $s^L$
& $r$
& $(L-1)s^L$
& $2(L-1)$
& $\floor{N/2}b$ \\
\bottomrule
\end{tabular}%
}
\end{table}

The Dragonfly bisection formula is the group-level non-strict bisection used
by the repository's topology-level estimate. For odd $g$, it partitions groups
as $\floor{g/2}$ and $\ceil{g/2}$. The Fat-tree bisection formula is the
terminal/full-bisection estimate for a regular full-bisection Fat-tree.

For UBMesh, the large bisection values come from the nD-FullMesh rule, not
from nearest-neighbor mesh links. Along dimension $i$, each line is a complete
graph $K_{L_i}$. Splitting that dimension into
$\floor{L_i/2}$ and $\ceil{L_i/2}$ routers cuts
$\floor{L_i/2}\ceil{L_i/2}$ links per line, and there are $R/L_i$ such lines.
For example, \toponame{ubmesh\_16x16\_c4} has $R=256$ routers and $16$
complete lines in either dimension, so the router-network bisection is
$16 \cdot 8 \cdot 8 \cdot 1\Gbps = 1024\Gbps$.
The concentration $c=4$ changes the terminal count and router radix, but it
does not change the router-router link count or router-network bisection
bandwidth.

\section{Concrete Examples Comparison}

Tables~\ref{tab:examples-256} and~\ref{tab:examples-1024} instantiate the
formulas with homogeneous links of latency $1$ cycle and bandwidth
$1\Gbps$. The ``No. nodes'' column counts terminal endpoints. The
\toponame{slimnoc} rows are not exactly 256 or 1024 nodes because the SlimNoC
finite-field construction only permits specific graph sizes.

\begin{table}[htbp]
\centering
\caption{Topology examples around 256 terminal nodes.}
\label{tab:examples-256}
\begin{tabular}{@{}lrrrrr@{}}
\toprule
Type & No. nodes & Router radix & No. links & Diameter & Bisection bandwidth \\
\midrule
\toponame{mesh2d\_16x16} & 256 & 5 & 480 & 30 & $16\Gbps$ \\
\toponame{mesh3d\_8x8x4} & 256 & 7 & 640 & 17 & $32\Gbps$ \\
\toponame{dragonfly\_p2\_a8\_h2} & 272 & 11 & 612 & 3 & $72\Gbps$ \\
\toponame{hypercube\_d8} & 256 & 9 & 1024 & 8 & $128\Gbps$ \\
\toponame{slimnoc\_q5\_p4} & 200 & 11 & 175 & 2 & $48\Gbps$ \\
\toponame{ubmesh\_8x8\_c4} & 256 & 18 & 448 & 2 & $128\Gbps$ \\
\toponame{ubmesh\_4x4x4\_c4} & 256 & 13 & 288 & 3 & $64\Gbps$ \\
\toponame{ubmesh\_4x4x2x2\_c4} & 256 & 12 & 256 & 4 & $32\Gbps$ \\
\toponame{ubmesh\_4x2x2x2x2\_c4} & 256 & 11 & 224 & 5 & $32\Gbps$ \\
\toponame{fattree\_r8\_l4} & 256 & 8 & 768 & 6 & $128\Gbps$ \\
\bottomrule
\end{tabular}
\end{table}

\begin{table}[htbp]
\centering
\caption{Topology examples around 1024 terminal nodes.}
\label{tab:examples-1024}
\begin{tabular}{@{}lrrrrr@{}}
\toprule
Type & No. nodes & Router radix & No. links & Diameter & Bisection bandwidth \\
\midrule
\toponame{mesh2d\_32x32} & 1024 & 5 & 1984 & 62 & $32\Gbps$ \\
\toponame{mesh3d\_8x8x16} & 1024 & 7 & 2752 & 29 & $64\Gbps$ \\
\toponame{dragonfly\_p4\_a8\_h4} & 1056 & 15 & 1452 & 3 & $272\Gbps$ \\
\toponame{hypercube\_d10} & 1024 & 11 & 5120 & 10 & $512\Gbps$ \\
\toponame{slimnoc\_q9\_p8} & 1296 & 21 & 1053 & 2 & $320\Gbps$ \\
\toponame{ubmesh\_16x16\_c4} & 1024 & 34 & 3840 & 2 & $1024\Gbps$ \\
\toponame{ubmesh\_8x8x4\_c4} & 1024 & 21 & 2176 & 3 & $256\Gbps$ \\
\toponame{ubmesh\_4x4x4x4\_c4} & 1024 & 16 & 1536 & 4 & $256\Gbps$ \\
\toponame{ubmesh\_4x4x4x2x2\_c4} & 1024 & 15 & 1408 & 5 & $128\Gbps$ \\
\toponame{fattree\_r8\_l5} & 1024 & 8 & 4096 & 8 & $512\Gbps$ \\
\bottomrule
\end{tabular}
\end{table}

\bibliographystyle{plain}
\bibliography{references}

\end{document}
